\documentclass[10pt]{scrartcl}
\usepackage[mdthm,colorsec,noasy]{evan}
\usepackage{handout}
\usepackage{figure-ink}
% Course-specific macros go here.
% Ink figure styles (as in the Berkeley handouts).
\newcommand{\PacketInkPath}[2][]{\InkStroke[#1]{#2}}
\tikzset{
  >={Stealth[length=4.6pt,width=2pt,inset=1.2pt]},
  figaxis/.style={ink arrow,draw=inkslate,line width=.35pt},
  figvec/.style={ink arrow,draw=inkblue},
  figvecsub/.style={ink arrow,draw=inkslate},
  figguide/.style={ink guide,draw=inkslate}
}
\title{Week N: TOPIC}
\subtitle{SOURCES (e.g.\ P\&B \S1.1--1.2; ACPS \S2.3; Engel Ch.~8)}
\author{Aaron Xiong}
\date{DATE}
\ihead{}
\begin{document}
\maketitle

\section{Facts used from earlier weeks}
% (T1), (T2), ...: \target{T1}\tool{Name (\xlink{week01}{anchor}{Week~1, Fact~2.5}).} statement, and the reason in one line.

\section{This week}
% For each topic, in lecture order:
%   \subsection{Topic (Ch.~X)}
%   \begin{qbox} what problem this section solves, with a small example \end{qbox}
%   For each result:
%     \begin{proposition*}[Notes number / textbook number] statement \end{proposition*}
%     \begin{ideabox} how to find the argument, in plain sentences or a small worked case (only when there is a real idea) \end{ideabox}
%     \begin{proof} \begin{steps} \ii ... \end{steps} \end{proof}
%   Figures with \begin{fig} ... \figcap{...} \end{fig}
%   \begin{shortcuts} ... \end{shortcuts}
%   \begin{trapbox} ... \end{trapbox}

\section{Problem set}
% Graded set (same problems as weekNN/pset.pdf): \begin{prob}\target{prob-1}\textbf{Problem 1 (P\&B 1).}\ statement\end{prob}
% Extra practice (real past problems): \begin{prob}\target{prob-X1}\textbf{X1 (Putnam 1999 A1).}\ statement\end{prob}

\section*{Hints for the extra practice}
% Graded problems get no printed hints (a hint on request costs points, syllabus Rules).

\end{document}
